\section{Study Design}
\label{sec:study}

\subsection{Research Questions}
\label{sec:study:rqs}

The study manipulates exactly one variable---how much of the shared workspace
the embedded agent can observe and act on (the \CI$\to$\CII$\to$\CIII\ context
ladder)---measured against each participant's own no-agent baseline (C0). The
questions stay inside what that manipulation supports: how much context an
embedded agent needs and how users work with it. We deliberately make
\emph{no} claim that an all-in-one workspace outperforms separate tools; this
design has no fragmented-tooling arm to test that, and consolidation serves
only as the apparatus and motivation (\S\ref{sec:related:fragmentation}), not
as a hypothesis.

\begin{description}
  \item[RQ1 (primary) --- Context breadth $\to$ effectiveness.]
    Does broadening the agent's context (chat $\to$ +board $\to$ +GitHub) make
    the user more effective, and on which task types?

  \item[RQ2 --- Reliance.]
    How much do users delegate to versus verify and override the agent, and
    does that division of labor shift with the agent's context breadth?

  \item[RQ3 --- Trust versus reality.]
    Does users' trust in the agent track its measured effectiveness, or diverge
    from it?
\end{description}

\subsection{Hypotheses}
\label{sec:study:hypotheses}

\begin{description}
  \item[H1~(RQ1, speed).]
    Broader agent context yields a larger improvement over the participant's own
    C0 baseline (faster completion): \CIII\ $>$ \CII\ $>$ \CI\ in
    baseline-normalized speed.

  \item[H2~(RQ1, task-type interaction).]
    The gain is \emph{task-dependent}: negligible on comprehension and
    single-surface tasks (where cross-surface data is irrelevant) but
    substantial on cross-surface tasks that span chat, board, and GitHub. This
    condition~$\times$~task-type interaction is the key result.

  \item[H3~(RQ1, cognitive load).]
    Broader context lowers self-reported cognitive load (\NASATLX) relative to
    baseline more than narrow context does, driven primarily by the effort and
    mental-demand subscales.

  \item[H4~(RQ2, reliance calibration).]
    Delegation to the agent rises with context breadth, but reliance is
    imperfectly calibrated: users do not verify or override precisely where the
    agent is actually wrong.

  \item[H5~(RQ3, perception--behavior gap).]
    Trust and perceived speedup need not track the agent's measured
    effectiveness: participants may trust an agent that \emph{does more} even
    when it is no faster, or is occasionally wrong. This perception--behavior
    divergence is the study's most novel angle.
\end{description}

\subsection{Design Overview}
\label{sec:study:design}

The study is a \textbf{hybrid within/between-subjects} design. Each participant
first completes six manual baseline tasks with no agent (the within-subjects C0
anchor), then is randomly assigned to one of three agent-context conditions and
completes six agent-assisted tasks in that fixed condition (the between-subjects
manipulation). The per-person C0 baseline lets each agent-task time be analyzed
as a \emph{delta from that participant's own manual speed}, absorbing the
individual skill variance that would otherwise sit between the between-subjects
cells and buying back much of the statistical power a pure between-subjects
design loses. Holding the agent's access fixed per person avoids within-person
contamination---once a participant has seen the full-context agent solve a task,
they cannot un-see how it should go---and keeps the session within a self-serve
$\sim$55--65-minute budget.

\subsection{Participants}
\label{sec:study:participants}

% TODO: Fill in once recruitment is complete.

We target $N$~=~20--30 participants drawn from two complementary populations:
\emph{university CS students} (undergraduate and graduate, recruited via course
mailing lists) and \emph{professional software developers} ($\geq$1 year of
industry experience, recruited via professional networks and online
communities). Recruiting both groups lets us examine whether the benefit of
context breadth---and reliance on and trust of the agent---varies with prior
experience. Eligibility further required self-reported familiarity with at least
one task-management tool (e.g., Jira, Linear, GitHub Projects) and one team-chat
tool (e.g., Slack, Teams, Discord). The study was conducted remotely.
Participants were compensated \emph{[amount TBD]}.

\subsection{Conditions}
\label{sec:study:conditions}

Participants were randomly assigned, with balanced permuted-block randomization
(block size 3), to one of three between-subjects conditions:

\begin{description}
  \item[\CI~(Chat only, advisory).] The agent observes only team chat history.
    No access to board state or GitHub. Advisory-only: no tools are offered to
    the model, so it can advise but cannot modify any surface.

  \item[\CII~(Chat + Board, action-capable).] The agent observes chat and the
    full board card state. Available tools: \texttt{create\_card},
    \texttt{move\_card}, \texttt{update\_card}, and \texttt{post\_message}
    (drafts a chat message for the participant to review before sending).

  \item[\CIII~(Chat + Board + GitHub, full action).] The agent observes all
    three surfaces. All \CII\ tools plus \texttt{link\_card\_to\_github},
    enabling it to connect board cards to GitHub issues or pull requests.
\end{description}

\paragraph{Capability scales with context---by design.} An agent cannot act on
a board it cannot see, so visibility and action capability are intentionally
coupled along the ladder. This bundling \emph{is} the phenomenon under study
(``does giving the agent more access help the user''); we state it as a design
decision, not a confound to be removed.

\subsection{Tasks}
\label{sec:study:tasks}

Each participant completes \textbf{12 tasks in fixed order}: six manual
\emph{baseline tasks} (B1--B6, C0, no agent) followed by six \emph{agent-assisted
tasks} (T1--T6) in their assigned condition. Each baseline is built to match an
agent task type, giving a like-for-like per-person normalizer, while baseline
and agent content are \emph{disjoint} to prevent learning transfer. The agent
phase uses the \texttt{gpt-4o-mini-2024-07-18} model at temperature 0.2 with
versioned, condition-specific system prompts pinned for reproducibility.

Tasks span three types across both phases:

\begin{description}
  \item[Comprehension.] Synthesize information from one or more surfaces to form
    a judgment (e.g., ``what does the team need most right now?''). All
    conditions should do well; chat-only is competitive.
  \item[Single-surface.] Completable by acting on a single surface; extra context
    should add little. These serve as controls showing broader context does
    \emph{not} universally help.
  \item[Cross-surface.] Explicitly requires coordinated action across two or more
    surfaces. Where H1/H2 live: full context should pull ahead, chat-only should
    struggle.
\end{description}

\begin{table}[h]
  \caption{Task design summary. B = Baseline (no agent, C0), T = Agent task.
    Surfaces: Ch = Chat, Bo = Board, GH = GitHub. T6 is constructed so the
    correct PR is identifiable \emph{only} with GitHub context---the single
    cleanest test of the top of the ladder.}
  \label{tab:tasks}
  \begin{tabular}{lllll}
    \toprule
    ID & Phase & Type & Surfaces & Summary \\
    \midrule
    B1 & Baseline & Single-surface & Ch, Bo     & Find blocked card; move to In Progress \\
    B2 & Baseline & Cross-surface  & Ch, Bo     & Check a card's column; post status reply \\
    B3 & Baseline & Cross-surface  & Ch, Bo     & Open card from chat link; set priority \\
    B4 & Baseline & Cross-surface  & GH, Ch     & Find open PR in GitHub view; post number in chat \\
    B5 & Baseline & Cross-surface  & Ch, Bo     & Recover a decision from chat; record on a card \\
    B6 & Baseline & Comprehension  & Bo, Ch     & Find CI-pipeline card; confirm what it does in chat \\
    \midrule
    T1 & Agent & Comprehension  & Ch, Bo     & ``What needs doing?'' then act \\
    T2 & Agent & Comprehension  & Ch, Bo, GH & ``Catch me up'': summarize recent activity \\
    T3 & Agent & Single-surface & Ch         & Draft and post a status update in chat \\
    T4 & Agent & Single-surface & Bo         & Set priorities on backlog cards \\
    T5 & Agent & Cross-surface  & Ch, Bo     & Take a bug from chat to the board \\
    T6 & Agent & Cross-surface  & Ch, Bo, GH & Link the PR to the card; note it in chat \\
    \bottomrule
  \end{tabular}
\end{table}

Because condition is between-subjects, all participants see the \emph{same}
tasks; only their agent's access differs. Task content is thus held constant and
context breadth is the sole manipulation.

\subsection{Procedure}
\label{sec:study:procedure}

% TODO: Fill in IRB protocol details, session timing.

Each session proceeded as follows:
\begin{enumerate}
  \item Informed consent and eligibility screening.
  \item Orientation to \Weld\ (un-timed, guided).
  \item \textbf{Baseline phase} (6 tasks, no agent, condition C0): tasks
    presented sequentially; participants self-paced and rated each task before
    advancing.
  \item \textbf{Post-baseline interstitial}: a \NASATLX\ instrument is
    administered here, \emph{before any agent exposure}
    (\texttt{tlx\_*\_baseline}), capturing manual-phase load. The agent panel is
    then revealed; \CII/\CIII\ participants receive a brief explanation of the
    available agent actions.
  \item \textbf{Agent phase} (6 tasks): same structure, with the agent available
    under the participant's assigned condition.
  \item Exit survey: a second \NASATLX\ (\texttt{tlx\_*}, agent phase), \SUS,
    integration-perception items, the AI-trust block, open-ended comments, and
    demographics.
  \item Debrief.
\end{enumerate}

Sessions were conducted remotely. Screen recording was \emph{not} collected; all
behavioral data came from in-product telemetry (\S\ref{sec:system:telemetry}).

\subsection{Measures}
\label{sec:study:measures}

\paragraph{Task performance (primary, RQ1).}
End-to-end \emph{active} time-on-task per task---the client clock runs only while
the tab is visible and persists across reloads, so idle and away time are
excluded by construction (\texttt{task\_completed.task\_time\_ms}). The headline
outcome is \textbf{baseline-normalized time}: each agent-task time divided by that
participant's own type-matched C0 baseline---comprehension (T1, T2)~$\div$~B6,
single-surface (T3, T4)~$\div$~B1, cross-surface (T5, T6)~$\div$~mean(B2--B5).
Completion is judged from \texttt{task\_final\_state}.

\paragraph{Correctness (RQ1, quality control).}
Per-task correctness on a 0--3 rubric, hand-coded by two raters (Cohen's $\kappa$
on a shared subset). Always reported beside time: a fast-but-wrong agent is not
effective.

\paragraph{Cognitive load (H3).}
\NASATLX~\cite{hart1988tlx}, six subscales, collected \textbf{twice}---
post-baseline (\texttt{tlx\_*\_baseline}) and at exit (\texttt{tlx\_*})---so load
is analyzed as a within-person agent$-$baseline delta rather than a single
end-of-session snapshot. A per-task effort rating (1--7, every task including
baselines) provides a finer-grained, task-localized companion measure.

\paragraph{Reliance and division of labor (RQ2).}
From telemetry: the \emph{delegation ratio} (agent- vs.\ user-performed actions),
verification behavior (re-reads, edits to the agent's output, undo), and
acceptance vs.\ override of agent suggestions
(\texttt{agent\_suggestion\_accepted/\_rejected}, \texttt{agent\_output\_edited},
\texttt{agent\_action\_undone}), plus view-switching and cross-reference usage.

\paragraph{Trust (RQ3).}
An eight-item AI-trust block (adapted from~\cite{madsen2000measuring}) covering
reliability, predictability, competence, willingness to rely without checking,
and whether the agent had enough context, plus \texttt{trust\_speed\_perceived}.
The gap between perceived speedup/trust and measured normalized time/correctness
operationalizes H5.

\paragraph{Usability and integration perception (descriptive).}
\SUS~\cite{brooke1996sus} (10 items) and a five-item integration-perception
scale; \texttt{int\_switch\_perceived} is paired with the behavioral view-switch
rate. These describe the apparatus and are not RQ outcomes.

\paragraph{Scoping manipulation check.}
The per-turn \texttt{agent\_context\_sent} snapshot (\S\ref{sec:system:gate})
provides direct evidence that the IV held: every \CI\ snapshot must contain no
board or GitHub context, and every \CII\ snapshot no GitHub context. Reported
before any inferential result.

\subsection{Analysis}
\label{sec:study:analysis}

The primary RQ1 model is a linear mixed model on log-transformed normalized
time: \texttt{log(norm\_time) \textasciitilde\ condition * task\_type +
(1\,|\,participant)}, with the condition~$\times$~task-type interaction as the
H2 test. Correctness is modeled in parallel with an ordinal cumulative-link
mixed model. H3 compares the within-person \NASATLX\ delta between conditions.
RQ2 is reported as reliance descriptives (delegation ratio, the
invoke\,$\to$\,execute\,$\to$\,accept/edit/undo funnel) by condition, with
thematic analysis of the open-ended items. RQ3 (H5) Spearman-correlates the
trust composite and \texttt{trust\_speed\_perceived} against measured normalized
time and correctness, characterizing divergence. Given three between-subjects
cells at $N\approx$~7--10 each, we lead with effect sizes and 95\% CIs over
null-hypothesis tests and pre-register the analysis.
